\section{AI Adoption：从 Mr. T Demo 到 Category-Defining Workflow}

访谈最后，McKinnon 用早期 iPhone 应用类比 AI：硬件/模型能力显然重要，但最初大量应用只是把新能力包装成 novelty，真正改变行业的是后来才出现的 Uber 式 workflow。这个判断不是说 AI 没价值，而是提醒 incumbent 和创业团队区分 capability shock、demo、repeatable workflow 与 organization redesign。

\subsection{技术能力与组织惯性}

大型公司拥有客户、数据和流程，也有 role、budget、support model 与文化惯性。即使 AI agent 能完成部分 support、operations 或 security task，组织仍需重写 accountability、quality review、exception handling、identity delegation 和 performance metric。若只在旧流程前加聊天框，系统可能得到一个“Mr. T app”，却没有改变 outcome economics。

\lecturefigure{16-ai-adoption-curve.png}{AI 通常先经历 capability shock 与 novelty app，随后才出现 category-defining workflow 和组织重设计。}{本地字幕 36:00--39:02；概念重绘。}

\begin{knowledgebox}{读图：判断 killer workflow 的四个问题}
它是否完成以前无法经济完成的任务，而非只换 UI？使用是否产生 compounding data/feedback？失败是否有明确 owner 和 recovery？组织是否愿意改变 role、process、control 与 pricing？如果答案都是否，应用可能仍是有趣 demo；如果 workflow、distribution 与 governance 同时变化，技术才可能形成新的 category。
\end{knowledgebox}

\teachervoice{McKinnon 不是 AI doom 叙事，也没有声称 killer use case 已经确定。他把当前阶段比作 iPhone 早期会说俏皮话的 Mr. T 应用：大家知道平台重要，但“像 Uber 一样因新能力而诞生的全新工作流”仍在形成。}

\begin{warningbox}{Incumbent inertia 既是风险也是安全缓冲}
旧流程会阻碍有价值的自动化，也可能暂时阻止未经治理的 agent 获得生产权限。组织重设计应同时移除无效 handoff 和增加 delegation、review、audit、fallback；“AI-first”不能成为绕过 identity、change management 与 incident response 的理由。
\end{warningbox}

\subsection{本章小结}

AI 价值从 capability 到 workflow，再到 organization。Identity infrastructure 在这条路径上不是后置登录组件，而是 agent accountability、least privilege 和 cross-tool audit 的前提。

